\documentclass[11pt]{article}
\usepackage[T1]{fontenc}
\usepackage{lmodern}
\usepackage[letterpaper,margin=1in]{geometry}
\usepackage{amsmath,amssymb}
\usepackage{array,booktabs}
\usepackage{enumitem}
\usepackage[most]{tcolorbox}
\usepackage{xcolor}
\usepackage[colorlinks=true,linkcolor=teal!60!black,urlcolor=teal!60!black,citecolor=teal!60!black]{hyperref}
\usepackage[backend=biber,style=numeric-comp,sorting=ynt,maxbibnames=8,maxcitenames=2,giveninits=true]{biblatex}
\addbibresource{references.bib}
\newtcolorbox{callout}{colback=black!3,colframe=black!25,boxrule=0.4pt,arc=2pt,
  left=6pt,right=6pt,top=4pt,bottom=4pt,breakable}
\setlist{topsep=3pt,itemsep=1pt,parsep=1pt,leftmargin=1.4em}
\setlength{\parindent}{0pt}\setlength{\parskip}{0.5em}
\sloppy
\begin{document}
\section*{Survey 01 --- The Landscape of Mechanistic Interpretability}

\begin{callout}
\textbf{Series note.} This is the first survey in \texttt{literature/survey/}. It is a \emph{map}: where the 22 archived papers sit, how they group by approach, and --- the part written for you specifically --- what mathematics they \textbf{share} versus where they \textbf{fundamentally diverge}. It is deliberately high-altitude; per-paper deep dives and a focused critique of mathematical blind spots are for later surveys (Part VI sketches the agenda).


\textbf{Audience \& conventions.} Written for a mathematician + systems engineer who is not an ML specialist. ML jargon is offloaded to the \emph{ML~Concepts} primer; the first use of a term defined there is \emph{italicised}. Paper links point into \texttt{../papers/}. All 22 papers were read in full for this survey.


\textbf{Notation.} For a linear map $W:A\to B$ we distinguish (per the \emph{ML~Concepts} primer's convention) the \textbf{dual map} $W^{\vee}:B^{*}\to A^{*}$ --- canonical, metric-free (the symbol is non-standard; some write ${}^{t}W$ or $W'$) --- from the \textbf{adjoint} $W^{*}:B\to A$, which exists only once a metric is chosen. $W^{*}$ always signals an invoked metric; $W^{\vee}$ never does.

\end{callout}


\subsection*{0. The one-paragraph orientation}

Mechanistic interpretability (MI) tries to reverse-engineer trained neural networks into human-understandable algorithms. The corpus you collected is almost entirely the \textbf{Anthropic / "Transformer Circuits" lineage} (16 of 22 papers) plus a few external landmarks (Redwood's IOI circuit; Nanda et al.'s grokking; Park--Choe--Veitch's linear- representation geometry). Strikingly, this body of work is \textbf{mathematically monolithic at its foundation and pluralistic at its methods}: nearly every paper stands on the same four linear-algebraic commitments (Part I), but they split into six methodological camps that operationalise those commitments very differently (Part II). The deepest intellectual action --- and the place a mathematician can contribute --- is at two seams: (a) the \textbf{inner-product / gauge} question that the field mostly leaves implicit (Part III), and (b) the \textbf{linear direction vs. nonlinear manifold} fork that has only just opened (Part IV).



\subsection*{Part I --- The shared mathematical substrate}

Four commitments recur in essentially every paper. They are not independent; they interlock into a single worldview. Understanding these four is most of understanding the field.


\subsubsection*{I.1 The Linear Representation Hypothesis (LRH): \emph{features are directions}}

The foundational modelling assumption: a high-level concept ("this text is in French", "the subject is a person") is encoded as a \textbf{direction} $d_i$ (a unit vector, sometimes a ray or cone to fix sign) in some activation space, and the activation vector decomposes as a sparse linear combination \[ x \;\approx\; b + \sum_i f_i(x)\, d_i, \qquad f_i(x)\ge 0,\ \text{few nonzero.}\] "Linear" means two things at once: features \textbf{superpose additively}, and a feature's \textbf{intensity} is the scalar coefficient $f_i$. This single hypothesis underwrites \textcite{toy-models-superposition}, \textcite{towards-monosemanticity}, \textcite{scaling-monosemanticity}, \textcite{linear-representation-hypothesis-geometry}, and the entire SAE program. It is stated as an explicit, falsifiable hypothesis --- and Part IV is about the first serious crack in it.


\subsubsection*{I.2 Superposition: \emph{$\gg d$ features in $d$ dimensions, paid for in interference}}

If features are directions and there are far more candidate features than dimensions, the model packs them as an \textbf{overcomplete, almost-orthogonal frame}. Two facts of different \emph{character} are in play, and separating them pays off later (Part III):

\begin{itemize}
\item \emph{Existence is metric-free.} That you can have more features than dimensions at all is just linear dependence. Writing the \textbf{embedding map} $W: e_i \mapsto d_i$ (feature-coordinate space $\to$ the width-$d$ activation space), \textcite{toy-models-superposition}'s crispest statement is: superposition is present \textbf{iff $W$ is non-injective} ($\ker W\neq 0$; the directions $\{d_i\}$ linearly dependent) --- a kernel-and-rank claim needing no inner product, forced the moment \#features $> d$. (The familiar form "$W^{*}W$ singular" says the same through the adjoint.)
\item \emph{The cost is metric-dependent.} What superposition \emph{pays} is \textbf{interference}: the off-diagonal of the Gram matrix $G = D^{*}D$ (built from the adjoint, hence metric-laden). That this stays small --- the frame is \emph{almost}-orthogonal --- is a claim about an inner product, underwritten by \emph{Johnson--Lindenstrauss} (exponentially many near-orthogonal unit vectors exist in $\mathbb{R}^d$) and \emph{compressed sensing} (sparse codes are recoverable from low-dimensional linear maps). \emph{Which} inner product is left implicit --- exactly the question Part III says the field rarely confronts.
\end{itemize}

Everything downstream is a negotiation with that Gram matrix.


\subsubsection*{I.3 The transformer is \emph{almost} linear: residual stream, path expansion, QK/OV}

\textcite{framework} is the keystone. Its move: the \emph{residual stream} is a purely \textbf{linear communication channel} (blocks add their outputs), so the forward pass can be \textbf{path-expanded} into a sum of end-to-end linear paths, with the \emph{softmax} attention pattern $A$ as the only nonlinearity. \textbf{Freeze $A$} and the model is an exact linear function of the input tokens. Two low-rank objects fall out of every attention head: \[ W_{QK} : V\to V^{*} \;\;(\text{a bilinear form — } \textit{where}\text{ to attend; coordinates } W_Q^{\top}W_K), \qquad    W_{OV} = W_O W_V : V\to V \;\;(\text{an endomorphism — } \textit{what}\text{ to write}).\] The type signature \emph{is} the where/what split: that $W_{QK}$ lands in the \textbf{dual} $V^{*}$ (a pairing of two stream vectors into a score) is exactly why a transpose appears in its coordinate form $W_Q^{\top}W_K$, while $W_{OV}$ --- a map of stream vectors to stream vectors --- needs none; see the \emph{ML~Concepts} primer §1.4 and the adjoint/dual convention there. "\emph{Virtual weights}" $W^{(2)}_{\text{in}} W^{(1)}_{\text{out}}$ connect any two blocks through the stream. This algebra --- tensor/Kronecker products, low-rank factorisations, eigenanalysis of the OV "copying" matrix --- is the literal skeleton of Groups 1 and 4 below, and reappears as the QK bilinear form in Group 5.


\subsubsection*{I.4 No privileged basis (a $GL(V)$ gauge symmetry) --- except where a nonlinearity breaks it}

Because the residual stream is read and written only by linear maps, you can conjugate the whole network by any invertible $M\in GL(V)$ --- replace each read map $W_{\text{in}}$ by $W_{\text{in}}M^{-1}$ and each write map $W_{\text{out}}$ by $M W_{\text{out}}$ --- and get a \textbf{bit-identical function}. The residual stream therefore has \textbf{no privileged basis}: its coordinates are individually meaningless. (Two things \emph{do} survive this gauge: a read map's \textbf{kernel} and a write map's \textbf{image} transform covariantly, $\ker\mapsto M\ker$, $\operatorname{im}\mapsto M\operatorname{im}$, so \emph{containments between them} --- and hence inter-block \textbf{virtual weights} $W_{\text{in}}^B W_{\text{out}}^A$ --- are invariant, whereas any complement, norm, or angle is not; see the \emph{ML~Concepts} primer §1.3.) A \emph{pointwise nonlinearity} (ReLU on MLP neurons) breaks the symmetry and singles out a \textbf{privileged basis} --- which is why neuron-level interpretability is even conceivable. \textcite{mech-interp-variables-bases} names this; \textcite{privileged-bases-residual-stream} shows the symmetry is \emph{empirically} broken in the residual stream too (outlier coordinates, kurtosis $\gg 3$) and fingers the \emph{Adam} optimiser's per-coordinate normalisation as the culprit.


\begin{callout}
\textbf{Why these four are really one idea.} I.1 says meaning lives in directions; I.4 says the ambient space has no canonical directions; I.2 says we cram more directions than dimensions; I.3 says the dynamics shuttling those directions around are (frozen-)linear. Put together: \emph{the model computes by moving an overcomplete frame of direction-encoded features through a gauge-symmetric linear channel punctuated by sparse nonlinearities.} Hold that sentence; the whole field is footnotes to it.

\end{callout}


\subsection*{Part II --- The six approaches}

\begingroup\footnotesize\noindent
\begin{tabular}{>{\raggedright\arraybackslash}p{0.03\textwidth}>{\raggedright\arraybackslash}p{0.14\textwidth}>{\raggedright\arraybackslash}p{0.18\textwidth}>{\raggedright\arraybackslash}p{0.30\textwidth}>{\raggedright\arraybackslash}p{0.29\textwidth}}
\toprule
\# & Approach (camp) & Core question & Papers & Signature mathematics \\
\midrule
1 & \textbf{Algebraic circuit analysis} & Rewrite the forward pass into readable terms & \textcite{framework}, \textcite{induction-heads}, \textcite{interpretability-in-the-wild-ioi} & tensor/Kronecker products, low-rank $W_{QK},W_{OV}$, path expansion, OV eigenvalues, causal patching \\
2 & \textbf{Superposition geometry} & How do $\gg d$ features pack into $d$ dims? & \textcite{toy-models-superposition}, \textcite{superposition-memorization-double-descent}, \textcite{distributed-representations-composition-superposition}, \textcite{privileged-bases-residual-stream}, \textcite{toy-model-of-interference-weights} & Thomson problem, uniform polytopes / tegum products, JL, compressed sensing, Gram matrix $D^{*}D$, phase transitions, kurtosis test \\
3 & \textbf{Dictionary learning (SAEs)} & \emph{Recover} the feature frame from activations & \textcite{towards-monosemanticity}, \textcite{scaling-monosemanticity}, \textcite{softmax-linear-units}, \textcite{sparse-crosscoders-model-diffing} & overcomplete dictionary learning, $\ell_1$/$\ell_0$ sparse coding, encoder/decoder frames, scaling laws \\
4 & \textbf{Circuit tracing / attribution} & Read off causal computation graphs on a real model & \textcite{circuit-tracing-computational-graphs}, \textcite{on-the-biology-of-a-large-language-model}, \textcite{tracing-attention-feature-interactions}, \textcite{progress-on-attention} & local linearization (Jacobian/Taylor), stop-gradients, Neumann series $(I-A)^{-1}$, transcoders, interference-corrected virtual weights \\
5 & \textbf{Representation geometry} & Take the \emph{geometry} (metric, curvature, group structure) seriously & \textcite{linear-representation-hypothesis-geometry}, \textcite{when-models-manipulate-manifolds-counting}, \textcite{progress-measures-for-grokking} & non-Euclidean inner products, Riesz duality, differential geometry of curves, Fourier analysis / characters of $\mathbb{Z}/n\mathbb{Z}$ \\
6 & \textbf{Epistemics / vision} & What makes an interpretation \emph{correct}? Where is this going? & \textcite{interpretability-dreams}, \textcite{toy-model-of-mechanistic-unfaithfulness}, \textcite{mech-interp-variables-bases} & faithfulness vs. function-equality, Jacobian matching, the program's axioms \\
\bottomrule
\end{tabular}
\endgroup


\subsubsection*{Group 1 --- Algebraic circuit analysis}
\emph{The exact-rewrite tradition.} \textcite{framework} treats the (attention-only) transformer as a tensor-algebraic object and multiplies it out; 0-layer models are bigrams $W_U W_E$, 1-layer models are skip-trigrams readable from $W_U W_{OV} W_E$ and $W_E^{\vee} W_{QK} W_E$ (the dual map $W_E^{\vee}$ pulls the QK form back to token space), and 2-layer models gain power through \textbf{composition} of heads, yielding the \textbf{induction head}. \textcite{induction-heads} elevates that single circuit to a theory of \emph{in-context learning} and ties its sudden formation to a \textbf{phase change} in training. \textcite{interpretability-in-the-wild-ioi} (the one Redwood paper) reverse-engineers a 26-head circuit in GPT-2 and contributes the field's first \textbf{formal validation criteria} --- \emph{faithfulness, completeness, minimality} --- plus \textbf{path patching}. Mathematically this group is the most classical: it is linear algebra and causal-intervention bookkeeping. Its open wound, stated in \textcite{framework} itself, is that \textbf{MLPs resist the rewrite} (the nonlinearity won't freeze) --- which is the cue for Group 3.


\subsubsection*{Group 2 --- Superposition geometry}
\emph{The forward-geometry tradition: given features and sparsity, what packing does training choose?} \textcite{toy-models-superposition} is the gem here: in controlled ReLU autoencoders it shows superposition is governed by a sparsity-vs-importance \textbf{phase transition}, and that the learned feature directions self-organise into \textbf{solutions of a generalized Thomson problem} --- antipodal pairs, pentagons, tetrahedra, square antiprisms --- often decomposing as \textbf{tegum (orthogonal) products} of low-dimensional uniform polytopes. The capacity question gets a compressed-sensing answer: \#features is linear in $d$ but modulated by sparsity. The companion notes specialise it: \textcite{superposition-memorization-double-descent} (finite data $\to$ memorised \emph{datapoints} live in superposition as polygon vertices; double descent falls out); \textcite{distributed-representations-composition-superposition} (separates \textbf{composition} from \textbf{superposition} as two ways to spend the same $\exp(d)$ representational volume --- and argues they are \emph{in tension}); \textcite{privileged-bases-residual-stream} (the symmetry-breaking diagnostic); and \textcite{toy-model-of-interference-weights} (the \emph{weights} between superposed features are themselves forced into superposition, producing spurious "interference weights" = noise). This group is where the most beautiful, most classical mathematics already lives.


\subsubsection*{Group 3 --- Dictionary learning (Sparse Autoencoders)}
\emph{The inverse problem to Group 2.} If the model stores an overcomplete sparse code, \textbf{recover the dictionary}. \textcite{towards-monosemanticity} trains \emph{SAEs} on a one-layer model's MLP and extracts thousands of monosemantic features; \textcite{scaling-monosemanticity} scales this to Claude 3 Sonnet (millions of features, scaling laws for SAE training, a sigmoid law relating a concept's training-frequency to whether it gets a dedicated feature). \textcite{softmax-linear-units} is the contrarian: instead of \emph{recovering} features post-hoc, \emph{change the architecture} ($\operatorname{SoLU}(x)=x\odot\operatorname{softmax}(x)$) so features align to neurons in the first place --- but a needed LayerNorm "smuggles superposition back", which the authors read as evidence \emph{for} the superposition hypothesis. \textcite{sparse-crosscoders-model-diffing} generalises SAEs to \textbf{crosscoders} (one dictionary across many layers/models). The whole group is \textbf{dictionary learning / sparse coding} with a learned overcomplete dictionary and an $\ell_1$ surrogate for $\ell_0$; the recurring anxiety is that there is \textbf{no ground-truth objective} --- reconstruction-MSE-plus-sparsity is only a \emph{proxy} for interpretability.


\subsubsection*{Group 4 --- Circuit tracing / attribution}
\emph{Scale Group 1's exact rewrite to real MLP-bearing models by replacing "exact" with "locally linear".} \textcite{circuit-tracing-computational-graphs} swaps each MLP for a \textbf{cross-layer transcoder} (a Group-3-style sparse feature model that predicts \emph{computation}, not just representation), freezes attention and norm denominators, adds per-token \textbf{error nodes} to match the model exactly on one prompt, and reads off a \textbf{linear attribution graph}. The math is \emph{local linearization}: stop-gradients make each feature's pre-activation an exact linear sum of upstream contributions (a Taylor expansion, exact at the prompt, divergent away from it); node influence is the \textbf{Neumann series} $(I-A)^{-1}-I$. \textcite{on-the-biology-of-a-large-language-model} applies it to Claude 3.5 Haiku across $\sim$10 case studies (planning, multi-hop reasoning, multilingual circuits, refusals). \textcite{tracing-attention-feature-interactions} and \textcite{progress-on-attention} attack the part the method \emph{freezes} --- the QK bilinear form $V\to V^{*}$ --- by decomposing the attention score into \textbf{feature--feature pairings} $\langle W_{QK}x_k, x_q\rangle$. This group is methodologically dominant today and is where "\emph{mechanistic faithfulness}" becomes the central worry (Group 6).


\subsubsection*{Group 5 --- Representation geometry}
\emph{The mathematically deepest group, and the one most aligned with your interests.} It treats the geometry as geometry rather than as bookkeeping. \textcite{linear-representation-hypothesis-geometry} formalises "feature = direction" via counterfactual concept pairs, distinguishes the \textbf{embedding} ($\Lambda$, input/context side) and \textbf{unembedding} ($\bar\Gamma$, output/word side) spaces --- which are \emph{canonically dual}, $\Lambda\cong\bar\Gamma^{*}$, the softmax pairing $\langle\lambda,\gamma\rangle$ needing no metric --- and --- crucially --- observes that the Euclidean inner product is \textbf{not canonical} (see Part III), introducing the \textbf{causal inner product} --- the metric $g_C=\operatorname{Cov}(\gamma)^{-1}:\bar\Gamma\xrightarrow{\sim}\bar\Gamma^{*}$, so $\langle\bar\gamma,\bar\gamma'\rangle_C=\langle g_C\,\bar\gamma,\,\bar\gamma'\rangle$ --- and a \textbf{Riesz-duality} theorem unifying the two spaces. \textcite{when-models-manipulate-manifolds-counting} finds that a scalar \emph{count} is represented not as a direction but as a \textbf{curved 1-D manifold} (a rippling helix in a $\sim$6-D subspace), and that an attention head \textbf{compares} two counts by a $W_{QK}$ that \textbf{rotates/twists one manifold onto another} --- differential geometry, not just linear algebra. \textcite{progress-measures-for-grokking} reverse-engineers modular addition as \textbf{Fourier multiplication}: inputs are placed on circles $(\cos\omega_k a,\sin\omega_k a)$ (the \textbf{characters of $\mathbb{Z}/n\mathbb{Z}$}) and addition becomes angle-addition, read out by trig identities --- representation theory of a finite abelian group. Note the resonance: \textbf{circular / periodic structure appears independently} in grokking (group characters) and in counting (the "ringing"/Gibbs geometry of the count manifold). That is a genuine shared mathematical phenomenon hiding under two different tasks.


\subsubsection*{Group 6 --- Epistemics / vision}
\emph{What does it mean to be right?} \textcite{mech-interp-variables-bases} sets the program's vocabulary (network $\approx$ compiled binary; features $\approx$ variables). \textcite{interpretability-dreams} is the manifesto: once superposition is solved, "for small enough circuits, statements about their behaviour become questions of \emph{mathematical reasoning}" --- an explicitly \textbf{axiomatic} ambition. \textcite{toy-model-of-mechanistic-unfaithfulness} is the cold shower: a transcoder can be monosemantic, sparse, and low-error \textbf{yet implement a different algorithm} that generalises differently off-distribution. Its proposed fix --- \textbf{Jacobian matching} (penalise $\|J_{\text{MLP}}-J_{\text{transcoder}}\|_F^2$) --- is mathematically the most pointed idea in the group, and it quietly concedes that matching input--output behaviour (Group 3's objective) is \emph{not} matching mechanism.



\subsection*{Part III --- What the groups SHARE (the connective tissue)}

This is the section you asked to be emphasised. Under the methodological diversity, the groups are bound by a small number of mathematical objects. Seeing the corpus this way is, I think, the single most useful reframe.


\subsubsection*{III.1 Almost everything is a bilinear pairing --- of two kinds: metric-free vs. metric-dependent}

Track the inner products through the corpus and nearly every quantity is a \textbf{bilinear pairing of two vectors}. The distinction that matters --- and that the field rarely flags --- is \emph{which kind}, because the two behave oppositely under the $GL(V)$ gauge of I.4.


\begin{itemize}
\item \textbf{Metric-free (gauge-invariant) pairings} are built only from the model's own maps, or from the natural contraction of a space with its dual; they need no choice of inner product:
\begin{itemize}
\item the \textbf{canonical pairing} $\langle\lambda(x),\gamma(y)\rangle$ the \emph{softmax} consumes (III.2) --- embedding contracted with unembedding, legitimate because the embedding space is the dual of the unembedding space, $\Lambda\cong\bar\Gamma^{*}$ (no metric needed);
\item the \textbf{attention score} $\langle W_{QK}x_k,\, x_q\rangle$ --- a \emph{learned} bilinear form $W_{QK}:V\to V^{*}$, invariant because its query- and key-sides transform contragrediently (\textcite{tracing-attention-feature-interactions} expands it into feature--feature terms);
\item a \textbf{linear probe} $\langle w,x\rangle$ --- a learned covector $w\in V^{*}$ evaluated on a vector.
\end{itemize}
\item \textbf{Metric-dependent pairings} pair two vectors \emph{of the same space through the standard inner product} --- equivalently they apply the \textbf{adjoint} ($W^{*}$, not the dual $W^{\vee}$; see the \emph{ML~Concepts} primer's convention), silently inserting $G=I$ as a metric the architecture never supplied:
\begin{itemize}
\item \textbf{interference / coherence} --- the off-diagonal of the feature Gram matrix $G=D^{*}D$ (Group 2); capacity, adversarial vulnerability, and the Thomson geometry are read off $G$;
\item \textbf{SAE reconstruction} $\|x-\hat x\|_2^2$ and the near-orthogonality $G\approx I$ that makes recovery work (Group 3); \emph{feature splitting} is $G$ revealing finer block structure;
\item \textbf{cosine similarity} between two feature directions (the default tool for clustering and comparing features).
\end{itemize}
\end{itemize}

The operational core of Groups 2--4 lives in the \emph{metric-dependent} column --- yet I.4 says the residual stream has no canonical metric. The one paper that treats the metric as something to \textbf{derive} rather than default to $I$ is the LRH paper, whose \textbf{causal inner product} answers "which $G$?" with $\operatorname{Cov}(\gamma)^{-1}$ (a whitening / Mahalanobis form). So:


\begin{callout}
\emph{Mathematically, the field is the study of one metric-free pairing (embedding$\cdot$unembedding, plus the learned form $W_{QK}$) alongside a host of metric-dependent Gram computations that quietly assume the very Euclidean structure the architecture declares arbitrary.}

\end{callout}

\subsubsection*{III.2 The gauge symmetry, and the quantity that actually survives it}

Combine I.4 (no privileged basis) with the LRH paper's identifiability law and you get a genuinely important --- and, in the corpus, almost entirely \textbf{unexploited} --- observation.


The softmax output depends only on the canonical pairing $\langle\lambda(x),\gamma(y)\rangle$, and is invariant under the simultaneous change (LRH paper, eq. 3.1) \[ \gamma(y)\ \mapsto\ A\,\gamma(y)+\beta, \qquad \lambda(x)\ \mapsto\ (A^{\vee})^{-1}\lambda(x), \qquad A\in GL(\bar\Gamma) \] (in coordinates $(A^{\vee})^{-1}=A^{-\top}$). This is the \textbf{readout instance of the I.4 gauge} ($A$ acting on the final-layer spaces): the unembedding side transforms by $A$, the embedding side by the inverse \textbf{dual map} $(A^{\vee})^{-1}$ --- \emph{contragrediently} --- so the pairing $\langle\lambda,\gamma\rangle$ is invariant (the additive $\beta$ merely shifts every logit by the same $y$-independent amount, which softmax discards --- §1.2). That $\lambda$ transforms as a covector is just the statement $\Lambda\cong\bar\Gamma^{*}$. The general principle, across every activation space:


\begin{callout}
\textbf{Gauge-invariant means built only from the model's own maps and the natural dual pairing.} Quantities assembled from learned weights and the contraction of a space with its dual --- the embedding--unembedding pairing $\langle\lambda,\gamma\rangle$, an attention score $\langle W_{QK}x_k, x_q\rangle$, a probe $\langle w,x\rangle$ --- are invariant. Quantities that pair two vectors of the \textbf{same} space through the \textbf{standard} inner product --- the cosine of two feature directions, $\|x-\hat x\|_2^2$, the Gram off-diagonal --- secretly insert the identity as a metric, and so are \textbf{gauge-dependent}: not canonically meaningful without first \emph{choosing} that metric.

\end{callout}

This has a sharp, slightly uncomfortable consequence for Group 3. An SAE minimises $\|x-\hat x\|_2^2$ in the residual stream and compares decoder directions by Euclidean cosine similarity --- both gauge-dependent quantities --- while Group 2/4 insist the same space has \emph{no privileged basis}. The practical escape hatch is empirical: training + Adam \emph{do} break the symmetry (\textcite{privileged-bases-residual-stream}), so the data-induced metric isn't arbitrary. But that is an empirical crutch, not a principled choice of metric --- and exactly one paper (the LRH paper) treats the metric as something to be \emph{derived} rather than assumed. \textbf{This is the central seam of the whole corpus} and I flag it as the leading candidate for mathematical contribution (Part VI).


The constructive flip side is worth stating, because the gauge problem is not a dead end: there \emph{is} a gauge-invariant way to ask whether two components interact --- not the cosine of their directions, but whether one's \textbf{image lands in the other's kernel}, equivalently whether their \textbf{virtual weight} $W_{\text{in}}^B W_{\text{out}}^A$ vanishes (I.4, the \emph{ML~Concepts} primer §1.3). Re-expressing interference and feature interaction in those metric-free terms --- rather than through the $G=I$ Gram matrix --- is one concrete shape the contribution in Part VI(1) could take.


\subsubsection*{III.3 The same low-rank attention algebra, reused five times}

$W_{QK}:V\to V^{*}$ (the bilinear form; coordinates $W_Q^\top W_K$) and $W_{OV}=W_O W_V:V\to V$ (the endomorphism) are introduced in \textcite{framework} and then \emph{never leave}: induction-head prefix-matching is a positive-eigenvalue OV plus a K-composition QK; the IOI circuit is classified entirely in these terms; circuit tracing freezes the QK and linearises the OV; the attention papers decompose QK into features; the manifold paper finds QK implementing a rotation. \textbf{If you learn one piece of transformer algebra, learn the QK/OV factorisation} --- it is the corpus's universal coordinate system.


\subsubsection*{III.4 Three regularizers, one Lagrangian shape}

Groups 2--4 and 6 all minimise \emph{reconstruction + sparsity}: $\ \mathbb{E}\|x-\hat x\|^2 + \lambda\,\rho(f)$, with $\rho$ an $\ell_1$ surrogate for $\ell_0$. Toy models, SAEs, transcoders, crosscoders, the interference-weight model, the unfaithfulness model --- same objective, different $x$, different decoder constraints. The shared pathologies (shrinkage from $\ell_1$, dead features, the $\ell_0$/$\ell_1$ gap exploited by "gradient masking") are therefore \emph{shared} across the whole tooling stack, not quirks of one method.


\subsubsection*{III.5 Compressed sensing as the two-way bridge}

Group 2 uses compressed sensing to \emph{upper-bound} how much superposition is possible (a forward statement about $W$); Group 3 \emph{is} the recovery algorithm compressed sensing promises (an inverse statement about recovering $f$ from $x=Wf$). Same theorem (RIP / incoherence / sparse recovery), read in both directions. This is the cleanest example of two camps sharing one piece of mathematics.



\subsection*{Part IV --- Where the math FUNDAMENTALLY diverges}

Genuine forks, not just different tools. These are the fault lines worth your attention.


\subsubsection*{IV.1 Feature-as-direction (linear) vs. feature-as-manifold (nonlinear) --- the big one}

Groups 1--4 assume a feature is a \textbf{1-D ray}: meaning is a direction, intensity is a scalar. \textcite{when-models-manipulate-manifolds-counting} breaks this for \emph{continuous} quantities: a count is a \textbf{curved 1-dimensional manifold} with deliberately high curvature ("ringing"), embedded in $\sim$6 dimensions, and the model exploits the curvature (a ray could only be scaled/compared by a monotone product, whereas a curve admits a genuine \textbf{rotation} that implements thresholded comparison). This is a \emph{category} difference: linear subspace vs. nonlinear submanifold; linear algebra vs. \textbf{differential geometry}. The question it raises is pointed and, as far as the corpus shows, open: \textbf{is the Linear Representation Hypothesis simply false for continuous/ordered variables, and right only for discrete categorical ones?} If so, the field needs a geometry of feature \emph{manifolds} (intrinsic dimension, curvature, how attention acts as a near-isometry between them) that currently does not exist in worked-out form. The LRH paper and the manifold paper are, in a precise sense, \textbf{mathematically incompatible accounts of what a feature is} --- and nobody has reconciled them. (Note the entanglement with III: the manifold's "curvature" and "ringing" are themselves read through the ambient inner product --- cosine-similarity matrices of activations --- so a properly metric-free account of feature \emph{manifolds} is doubly open.)


\subsubsection*{IV.2 Modelling representation vs. modelling computation}

SAEs (Group 3) model the \textbf{representation} --- a basis for activation \emph{space}. Transcoders and attribution graphs (Group 4) model the \textbf{computation} --- the map between activations. The \textcite{toy-model-of-mechanistic-unfaithfulness} paper shows these are \textbf{not interchangeable}: an SAE "can't really mislead you" (it just describes a vector), but a transcoder that reproduces the input--output map can still implement the \emph{wrong algorithm}. The mathematical content of the distinction: representation-modelling constrains a single space; computation-modelling must match a \textbf{Jacobian} (a linear map between spaces), and matching a function is strictly weaker than matching its derivative everywhere. This is a real divergence in what "correct" even means, and it cuts the toolkit in half.


\subsubsection*{IV.3 Exact algebra vs. local linearization vs. learned approximation}

Three incompatible notions of "linear" coexist (see the \emph{ML~Concepts} primer §1.4 and §3):

\begin{enumerate}
\item \textbf{Exact, by freezing the nonlinearity} (\textcite{framework}): attention frozen $\Rightarrow$ the forward pass is \emph{honestly} linear in tokens. No approximation.
\item \textbf{Local linearization} (\textcite{circuit-tracing-computational-graphs}, attribution patching): a Taylor/Jacobian expansion, \emph{exact at one prompt, wrong away from it.}
\item \textbf{Learned sparse approximation} (SAEs/transcoders): a trained model that is neither exact nor a derivative, only a low-error fit --- and possibly mechanistically unfaithful.
\end{enumerate}
These are routinely all called "linear", but they have different error semantics, and conflating them is where overclaiming creeps in. A mathematician will want to keep them strictly separate.


\subsubsection*{IV.4 Discrete-algebraic vs. continuous-geometric task structure}

\textcite{progress-measures-for-grokking} finds \textbf{finite group representation theory} ($\mathbb{Z}/n\mathbb{Z}$ characters); the manifold paper finds \textbf{continuous differential geometry}; the superposition papers find \textbf{discrete polytopes} (Thomson). These are three different branches of mathematics describing "the geometry of representations", and the field has \textbf{not} unified them. The tantalising hint (III, and the grokking$\leftrightarrow$counting circle resonance) is that \textbf{harmonic analysis / periodicity} might be the common generalisation --- Fourier on a group, Fourier features on a manifold, and the eigen- structure of a circulant Gram matrix are not unrelated --- but this is conjecture, not established.



\subsection*{Part V --- A few honest observations about the corpus}

\begin{itemize}
\item \textbf{Provenance bias.} 16/22 papers are one lab's (Anthropic) thread; they share authors, assumptions, and a house style (many are explicitly "preliminary lab-meeting notes"). The shared substrate of Part I is partly a sociological fact, not only a mathematical necessity. The three external papers (IOI, grokking, LRH-geometry) are also the three that contribute the most \emph{formal} machinery (validation metrics; Fourier reverse-engineering; the metric question) --- worth noting.
\item \textbf{Rigour gradient.} Formality decreases as models scale: toy models get closed-form loss analysis and named polytopes; frontier-model papers get qualitative case studies with $\sim$25\% success rates and heavy hedging. The mathematically tractable results are mostly in Group 2 (toy models); the high-stakes claims are mostly in Group 4 (frontier circuits), where faithfulness is weakest.
\item \textbf{Proxy objectives everywhere.} No one has a ground-truth definition of "feature" or a principled SAE objective; MSE+$\ell_1$ is admitted to be a stand-in. Several "definitions" (copying matrix via positive eigenvalues; interference weights; faithfulness) are explicitly heuristic and acknowledged as not-yet-formalisable.
\end{itemize}


\subsection*{Part VI --- Agenda for later surveys (where the math might be)}

Signposts, not conclusions. Each is a candidate "mathematical blind spot $\to$ contribution" to develop in a focused follow-up.


\begin{enumerate}
\item \textbf{A gauge theory of representations (III.2).} Make the $GL(V)$ symmetry explicit; characterise the invariants (the embedding--unembedding pairing, learned forms like $W_{QK}$, and image-in-kernel / virtual-weight relations); re-express superposition/SAE/probing results in gauge-invariant form; ask what an SAE \emph{should} minimise if not the gauge-dependent $\ell_2$ loss. The LRH paper's $\operatorname{Cov}(\gamma)^{-1}$ is one answer for the output space; the residual-stream interior is wide open.
\item \textbf{A geometry of feature manifolds (IV.1).} Reconcile direction-features with manifold- features. What is the right invariant description of a count/colour/date manifold (intrinsic dimension, curvature, the action of $W_{QK}$ as a near-isometry)? Is LRH a categorical- variable special case of a manifold theory?
\item \textbf{The interference Gram matrix as the central object (III.1).} A unified treatment of $G=D^{*}D$ --- spectrum, block structure (feature splitting), the relation between its off-diagonal (interference) and downstream "interference weights", and bounds linking sparsity, dimension, and the operator norm of $G-I$.
\item \textbf{Error semantics of the three "linearities" (IV.3).} Make precise when local linearization / learned approximation incurs how much error, and connect Jacobian-matching (mechanistic faithfulness) to a quantitative bound.
\item \textbf{Harmonic analysis as a possible unifier (IV.4).} Is there a single framework (Fourier on groups / spectral geometry / circulant structure) underlying the polytope, character, and manifold geometries? Likely the most speculative, possibly the most fruitful.
\end{enumerate}


\subsubsection*{Appendix --- paper $\to$ group quick index}

\begin{itemize}
\item \textbf{framework} --- G1 (substrate for all)
\item \textbf{induction-heads} --- G1
\item \textbf{interpretability-in-the-wild-ioi} --- G1 (+ validation metrics)
\item \textbf{toy-models-superposition} --- G2 (substrate for G3)
\item \textbf{superposition-memorization-double-descent} --- G2
\item \textbf{distributed-representations-composition-superposition} --- G2
\item \textbf{privileged-bases-residual-stream} --- G2 (substrate I.4)
\item \textbf{toy-model-of-interference-weights} --- G2 (bridges to G4)
\item \textbf{towards-monosemanticity} --- G3
\item \textbf{scaling-monosemanticity} --- G3
\item \textbf{softmax-linear-units} --- G3 (architectural variant)
\item \textbf{sparse-crosscoders-model-diffing} --- G3 (bridges to G4)
\item \textbf{circuit-tracing-computational-graphs} --- G4 (method)
\item \textbf{on-the-biology-of-a-large-language-model} --- G4 (application)
\item \textbf{tracing-attention-feature-interactions} --- G4/G1 (QK)
\item \textbf{progress-on-attention} --- G4/G1 (QK)
\item \textbf{linear-representation-hypothesis-geometry} --- G5 (the metric question)
\item \textbf{when-models-manipulate-manifolds-counting} --- G5 (the manifold fork)
\item \textbf{progress-measures-for-grokking} --- G5 (group/Fourier)
\item \textbf{interpretability-dreams} --- G6
\item \textbf{mech-interp-variables-bases} --- G6 (substrate I.4)
\item \textbf{toy-model-of-mechanistic-unfaithfulness} --- G6 (faithfulness)
\end{itemize}


\printbibliography
\end{document}
